% GENERATED FILE. Edit canonical lessons, then run npm run generate:books.
% canonical-source-hash: fnv1a64:9c813274f45e8fb9
% canonical-lessons: LA-C319-sipho, LA-C319-hama, LA-C319-restinguo, LA-C319-excubiae, LA-C319-amburo

\chapter{Vigilēs}
\label{ch:la-vigiles}
\hlchaptermodality{319}

\begin{chapteropening}
\textbf{By the end of this chapter:} \emph{I can tell how a fire was fought: the fire pump, the fire bucket, to put out, the night watch and to scorch.}

Complete the last lesson of chapter 319: I can tell how a fire was fought: the fire pump, the fire bucket, to put out, the night watch and to scorch.
\end{chapteropening}

\section[sīphō, sīphōnis]{sīphō, sīphōnis — a fire pump; a siphon}

\label{lesson:LA-C319-sipho}



\begin{quote}
Before the new one: say the Latin for to smoke, then the Latin for to go up in flames.
\end{quote}

\subsection*{You'll want to know: sīphō, sīphōnis}
\textbf{sīphō, sīphōnis} — \textquotedblleft{}a fire pump; a siphon\textquotedblright{}.

\begin{quote}
Oppidum nūllum sīphōnem habēbat.
\end{quote}

\begin{quote}
\emph{The town had no fire pump.}
\end{quote}

Masculine, 3rd declension, a Greek word. A Roman governor once wrote that a burning town had no \emph{sīphō} and no buckets.

\begin{cousinweb}
\textbf{sīphō} was borrowed from Greek \emph{siphōn}, \textquotedblleft{}a pipe\textquotedblright{} (Lewis \& Short): a small pipe, a drinking tube, and a fire engine. English \emph{siphon} comes from the same Greek word.
\end{cousinweb}

\begin{grammarlens}[title={What you've built}]
One more word for telling, explaining and describing at B1.
\end{grammarlens}

\subsection*{Your turn}
\begin{itemize}
\raggedright
  \item \emph{Say it:} \emph{sīphō, sīphōnis}
  \item \emph{Say it:} \emph{sīphō, sīphōnis}, once more
  \item \emph{Say it:} a sentence of your own with \emph{sīphō, sīphōnis}
  \item \emph{Recall:} say the Latin for to smoke, then the Latin for to go up in flames, then say \emph{sīphō, sīphōnis} again
\end{itemize}

\begin{tcolorbox}[breakable,colback=teal!4,colframe=teal!35!black,title={Before you move on}]
What does sīphō, sīphōnis mean? (\textquotedblleft{}A fire pump; a siphon\textquotedblright{}.) Say it once more.
\end{tcolorbox}
\section[hama, hamae]{hama, hamae — a fire bucket}

\label{lesson:LA-C319-hama}



\begin{quote}
Before the new one: say the Latin for to go up in flames, then the Latin for a fire pump.
\end{quote}

\subsection*{You'll want to know: hama, hamae}
\textbf{hama, hamae} — \textquotedblleft{}a fire bucket\textquotedblright{}.

\begin{quote}
Vigilēs aquam hamīs portābant.
\end{quote}

\begin{quote}
\emph{The watchmen were carrying water in fire buckets.}
\end{quote}

Feminine, short \emph{a}. The watchmen of the fire service kept them ready.

\begin{cousinweb}
\textbf{hama} was borrowed from Greek \emph{amē}, \textquotedblleft{}a bucket\textquotedblright{} (EtymDB, Lewis \& Short).
\end{cousinweb}

\begin{grammarlens}[title={What you've built}]
One more word for telling, explaining and describing at B1.
\end{grammarlens}

\subsection*{Your turn}
\begin{itemize}
\raggedright
  \item \emph{Say it:} \emph{hama, hamae}
  \item \emph{Say it:} \emph{hama, hamae}, once more
  \item \emph{Say it:} a sentence of your own with \emph{hama, hamae}
  \item \emph{Recall:} say the Latin for to go up in flames, then the Latin for a fire pump, then say \emph{hama, hamae} again
\end{itemize}

\begin{tcolorbox}[breakable,colback=teal!4,colframe=teal!35!black,title={Before you move on}]
What does hama, hamae mean? (\textquotedblleft{}A fire bucket\textquotedblright{}.) Say it once more.
\end{tcolorbox}
\section[restinguō, restinguere]{restinguō, restinguere — to put out, to quench}

\label{lesson:LA-C319-restinguo}



\begin{quote}
Before the new one: say the Latin for a fire pump, then the Latin for a fire bucket.
\end{quote}

\subsection*{You'll want to know: restinguō, restinguere}
\textbf{restinguō, restinguere} — \textquotedblleft{}to put out, to quench\textquotedblright{}.

\begin{quote}
Vigilēs incendium aquā restinxērunt.
\end{quote}

\begin{quote}
\emph{The watchmen put out the fire with water.}
\end{quote}

\emph{restinguō, restinguere, restinxī, restinctum}: close to \emph{exstinguō} (taught). Also of thirst: \emph{sitim restinguere}.

\begin{cousinweb}
\textbf{restinguō} is built on an old verb \emph{stinguō}, \textquotedblleft{}to quench\textquotedblright{}, with \emph{re-} (Lewis \& Short); \emph{exstinguō} is built on the same verb.
\end{cousinweb}

\begin{grammarlens}[title={What you've built}]
One more word for telling, explaining and describing at B1.
\end{grammarlens}

\subsection*{Your turn}
\begin{itemize}
\raggedright
  \item \emph{Say it:} \emph{restinguō, restinguere}
  \item \emph{Say it:} \emph{restinguō, restinguere}, once more
  \item \emph{Say it:} a sentence of your own with \emph{restinguō, restinguere}
  \item \emph{Recall:} say the Latin for a fire pump, then the Latin for a fire bucket, then say \emph{restinguō, restinguere} again
\end{itemize}

\begin{tcolorbox}[breakable,colback=teal!4,colframe=teal!35!black,title={Before you move on}]
What does restinguō, restinguere mean? (\textquotedblleft{}To put out, to quench\textquotedblright{}.) Say it once more.
\end{tcolorbox}
\section[excubiae, excubiārum]{excubiae, excubiārum — the watch, guard duty}

\label{lesson:LA-C319-excubiae}



\begin{quote}
Before the new one: say the Latin for a fire bucket, then the Latin for to put out.
\end{quote}

\subsection*{You'll want to know: excubiae, excubiārum}
\textbf{excubiae, excubiārum} — \textquotedblleft{}the watch, guard duty\textquotedblright{}.

\begin{quote}
Vigilēs per tōtam noctem excubiās agēbant.
\end{quote}

\begin{quote}
\emph{The watchmen kept watch all night long.}
\end{quote}

Feminine, plural only: lying out of doors on guard. \emph{Excubiās agere}: to keep watch.

\begin{cousinweb}
\textbf{excubiae} is built on \emph{ex}, \textquotedblleft{}out\textquotedblright{}, and \emph{cubō}, \textquotedblleft{}to lie down\textquotedblright{} (Lewis \& Short): a lying out on guard.
\end{cousinweb}

\begin{grammarlens}[title={What you've built}]
One more word for telling, explaining and describing at B1.
\end{grammarlens}

\subsection*{Your turn}
\begin{itemize}
\raggedright
  \item \emph{Say it:} \emph{excubiae, excubiārum}
  \item \emph{Say it:} \emph{excubiae, excubiārum}, once more
  \item \emph{Say it:} a sentence of your own with \emph{excubiae, excubiārum}
  \item \emph{Recall:} say the Latin for a fire bucket, then the Latin for to put out, then say \emph{excubiae, excubiārum} again
\end{itemize}

\begin{tcolorbox}[breakable,colback=teal!4,colframe=teal!35!black,title={Before you move on}]
What does excubiae, excubiārum mean? (\textquotedblleft{}The watch, guard duty\textquotedblright{}.) Say it once more.
\end{tcolorbox}
\section[ambūrō, ambūrere]{ambūrō, ambūrere — to scorch, to singe}

\label{lesson:LA-C319-amburo}



\begin{quote}
Before the new one: say the Latin for to put out, then the Latin for the watch.
\end{quote}

\subsection*{You'll want to know: ambūrō, ambūrere}
\textbf{ambūrō, ambūrere} — \textquotedblleft{}to scorch, to singe\textquotedblright{}.

\begin{quote}
Ignis manūs eius ambussit, sed vīvus fūgit.
\end{quote}

\begin{quote}
\emph{The fire scorched his hands, but he got away alive.}
\end{quote}

\emph{ambūrō, ambūrere, ambussī, ambustum}: to burn all round, not right through.

\begin{cousinweb}
\textbf{ambūrō} is built on \emph{ūrō}, \textquotedblleft{}to burn\textquotedblright{}, with \emph{amb-}, \textquotedblleft{}around\textquotedblright{} (EtymDB): to burn round the edges.
\end{cousinweb}

\begin{grammarlens}[title={What you've built}]
One more word for telling, explaining and describing at B1.
\end{grammarlens}

\subsection*{Your turn}
\begin{itemize}
\raggedright
  \item \emph{Say it:} \emph{ambūrō, ambūrere}
  \item \emph{Say it:} \emph{ambūrō, ambūrere}, once more
  \item \emph{Say it:} a sentence of your own with \emph{ambūrō, ambūrere}
  \item \emph{Recall:} say the Latin for to put out, then the Latin for the watch, then say \emph{ambūrō, ambūrere} again
\end{itemize}

\begin{tcolorbox}[breakable,colback=teal!4,colframe=teal!35!black,title={Before you move on}]
What does ambūrō, ambūrere mean? (\textquotedblleft{}To scorch, to singe\textquotedblright{}.) Say it once more.
\end{tcolorbox}
